%% ============================================================================
%% ============================================================================
\section{DAI's Documented USDC Exposure}
\label{app:dai}

DAI broke (\$0.886) through a contemporaneously documented exposure to USDC. The break was not a
separate coordination run, and counting it as one would double-count the USDC run, so this
appendix does not decompose DAI's $1{,}141$-bps de-peg into mechanical and panic components.

DAI was linked to USDC through MakerDAO's peg-stability module (PSM), which swaps
USDC$\leftrightarrow$DAI one-for-one up to a debt ceiling, and through USDC held as vault
collateral. A press report published at 10:10~UTC on 11~March 2023 and web-archived 44 minutes
later gives two distinct ratios in a single passage~\cite{theblock2023psm}. DAI was ``about 36\%
backed by USDC, according to daistats'' and the USDC PSM had ``reached its cap of 3.1 billion
USDC, so no more DAI can be minted with USDC''. Against the on-chain DAI supply at
2023-03-11 12:00~UTC (\$4.19B), the \$3.1B cap corresponds to $74\%$ of DAI outstanding having
been minted through the USDC PSM. The $74\%$ is a different quantity from the source's own
$36\%$ figure, which is a share of total collateral rather than of DAI supply. We report both
ratios (Table~\ref{tab:dai}) and use neither as a pass-through coefficient, because the two denominators are not
interchangeable.

Three facts prevent a decomposition. First, the two ratios differ by roughly a factor of two
and imply mechanical components of $444$ and $912$~bps against the same observed $1{,}141$~bps,
so any residual would be an artifact of which ratio we choose rather than a measurement.
Second, the source sentence that reports the \$3.1B also records that the USDC$\to$DAI mint
leg had reached its cap at the moment we analyze. That leg's arbitrage would
move DAI's price toward USDC's, so we cannot argue that the PSM raises effective pass-through
toward $100\%$. Consistent with the cap being binding, DAI traded at or above USDC through the
de-peg rather than converging to it. Third, we have no on-chain custody reading to compare
with the reported figure. The PSM's USDC had moved to off-chain custody by 2023, so the
module and its collateral adapter both hold approximately zero USDC at the run-date block,
and the \$3.1B stands as a reported debt-ceiling figure.

DAI carried a documented USDC exposure that USDT did not. DAI broke and USDT did not. DAI's
break is therefore not independent evidence of a second coordination run. We make no claim
about how much of DAI's $1{,}141$~bps was mechanical.

\begin{table}[t]
\centering
\caption{DAI's de-peg and its documented USDC exposure. We report both ratios the single
contemporaneous source gives, $74\%$ of DAI supply minted through the USDC PSM at its
\$3.1B cap and the source's own $36\%$ USDC share of DAI collateral, because they have
different denominators and imply materially different pass-through. We do not select
between them and report no mechanical/panic split. The \$3.1B is a cited reporting figure
(The Block, 11 March 2023) with no on-chain corroboration, for the reason given in
Appendix~\ref{app:dai}.}
\label{tab:dai}
\small
\checkwidth{\small% Source: results/number4_dai_contagion.json --
%   usdc_depeg_bps (1233.3 -> 1233); observed_dai_depeg_bps (1140.6 -> 1141);
%   dai_total_supply_usd (4193955100.3838315 -> $4.19B); reported_psm_usdc (3100000000.0 -> $3.1B);
%   psm_share_of_dai_supply (0.7391590815353001 -> 74%); source_collateral_share (0.36 -> 36%).
%   NO decomposition row: see the no_decomposition key in the same results file.
%   D19: the on_chain_psm_usdc_total row is retracted, not read here -- see
%   write_tab_dai docstring.
\begin{tabular}{@{}p{6.4cm}p{4.0cm}@{}}
\toprule
Quantity & Value \\
\midrule
USDC de-peg (trough, composite) & $1233$~bps \\
Observed DAI de-peg (trough, composite) & $1141$~bps \\
On-chain DAI supply, run-date block & \$4.19B \\
Reported PSM-USDC-A cap, 11 Mar 2023 & \$3.1B \\
\quad as a share of DAI supply & $74\%$ \\
Reported USDC share of DAI collateral & $36\%$ \\
\bottomrule
\end{tabular}
}
% Source: results/number4_dai_contagion.json --
%   usdc_depeg_bps (1233.3 -> 1233); observed_dai_depeg_bps (1140.6 -> 1141);
%   dai_total_supply_usd (4193955100.3838315 -> $4.19B); reported_psm_usdc (3100000000.0 -> $3.1B);
%   psm_share_of_dai_supply (0.7391590815353001 -> 74%); source_collateral_share (0.36 -> 36%).
%   NO decomposition row: see the no_decomposition key in the same results file.
%   D19: the on_chain_psm_usdc_total row is retracted, not read here -- see
%   write_tab_dai docstring.
\begin{tabular}{@{}p{6.4cm}p{4.0cm}@{}}
\toprule
Quantity & Value \\
\midrule
USDC de-peg (trough, composite) & $1233$~bps \\
Observed DAI de-peg (trough, composite) & $1141$~bps \\
On-chain DAI supply, run-date block & \$4.19B \\
Reported PSM-USDC-A cap, 11 Mar 2023 & \$3.1B \\
\quad as a share of DAI supply & $74\%$ \\
Reported USDC share of DAI collateral & $36\%$ \\
\bottomrule
\end{tabular}

\end{table}
